% Ad Familiares 2.15 --- headnote
% ad-familiares-02-15, 3 or 4 August 50 BC, Side

\textit{Cicero to M. Caelius Rufus, curule aedile,
written from Side on the return route, on the third or
fourth day before the Nones of Sextilis (3 or 4 August)
50 BC (Perseus dateline: \textit{Scr. Sidae iii aut
prid. Non. Sext. a. 704 (50)}). The Cilician term is
over: Cicero has handed off the province and is making
his way home along the Pamphylian coast, waiting on
the season's winds for the Aegean crossing. He is
answering a Caelius newsletter that has caught up with
him on the road, with its usual mixture of city news,
political weather, and prosecutorial gossip.}

\textit{Four short movements. Caelius has worked
deftly with Curio to extract the senatorial
supplication that crowns Cicero's modest campaign in
the Amanus; Cicero acknowledges the manoeuvre and
points forward to the triumph he now hopes will follow.
Then the family note --- Tullia's new husband Dolabella,
of whose marriage Cicero is privately uncertain, is
being praised and even befriended by Caelius, and
Cicero thanks him for the office of moderation. The
third section is the most pointed line of the letter:
``I favour Curio, I want Caesar to come off honourably,
for Pompey I am willing to die'' --- a triangulation
that already names the impossibility of the coming
year. The closing paragraph defends his choice to leave
the province in the hands of the quaestor Coelius
rather than his brother Quintus, with the
etesian-wind sign-off: unless the season pins him
down, he will see Caelius very soon.}
